% !TeX root = ../main.tex
\documentclass[../main]{subfiles}
\begin{document}
\chapter{Idea de Proyecto}
\img{images/idea.jpg}{La idea de proyecto es el motor que nos
  acompaña desde el inicio del proceso creativo. Conocemos el origen
  —la idea— y visualizamos el final: la materialización física y
funcional de nuestra propuesta.}
\info{Contenido}{
  \begin{itemize}
    \item Definición y alcance de la idea de proyecto.
    \item Diseño de la solución base optimizada.
    \item Identificación de necesidades técnicas y sociales.
    \item Organización y jerarquización de soluciones alternativas.
  \end{itemize}
}
\info{Solución Base Optimizada}{
  Una solución base optimizada es aquella diseñada específicamente
  para satisfacer los requisitos del proyecto con la máxima
  eficiencia. Esto requiere un análisis exhaustivo de requerimientos,
  restricciones técnicas y tecnologías disponibles.
  Para optimizar una propuesta en el campo electromecánico, consideramos:
  \begin{itemize}
    \item \textbf{Rendimiento y Escalabilidad:} Maximizar el sistema
      pensando en futuros crecimientos.
    \item \textbf{Seguridad y Mantenimiento:} Proteger
      usuarios/equipos y facilitar actualizaciones para minimizar
      tiempos de inactividad.
    \item \textbf{Costo y Sostenibilidad:} Minimizar el costo total
      de propiedad a largo plazo con una solución económica y duradera.
  \end{itemize}
}
\subsection{Conceptualización de la Idea}
La idea de proyecto es la propuesta inicial para resolver un problema
técnico o aprovechar una oportunidad de mejora. Es la fase
embrionaria donde se definen los pilares que luego serán
desarrollados con rigor técnico.

Toda idea de proyecto profesional debe estructurarse bajo los siguientes puntos:
\begin{description}
  \item[Diagnóstico:] Descripción precisa del problema u oportunidad a resolver.
  \item[Metas Técnicas:] Objetivos específicos que se pretenden alcanzar.
  \item[Alcance y Límites:] Definición clara de qué incluye el
    proyecto y qué queda fuera de él.
  \item[Gestión de Recursos:] Determinación de presupuesto, personal,
    materiales y tiempos.
  \item[Plan de Ejecución:] Hoja de ruta detallada con plazos y pasos a seguir.
\end{description}

\actividad{Laboratorio de Ideas}

\begin{enumerate}
  \item \textbf{Organización:} Establecer los grupos de trabajo para
    el desarrollo del proyecto anual.
  \item \textbf{Investigación y Diagnóstico:}
    \begin{enumerate}
      \item Realizar relevamientos o encuestas para detectar fallas o
        necesidades en la comunidad o el taller.
      \item Investigar tendencias tecnológicas y problemáticas
        sociales actuales mediante fuentes técnicas.
      \item Listar y priorizar los problemas encontrados según su
        viabilidad técnica.
    \end{enumerate}
  \item \textbf{Desarrollo de Soluciones:}
    \begin{enumerate}
      \item Proponer diferentes alternativas para los problemas detectados.
      \item Definir el grupo de características técnicas o límites
        (normativas) que deben cumplir las soluciones.
      \item Utilizar una matriz de ponderación para seleccionar la mejor opción.
      \item Definir la \textbf{Solución Base Optimizada} y proponer
        al menos 4 alternativas técnicas secundarias.
    \end{enumerate}
  \item \textbf{Formalización:} Redactar la idea de proyecto
    definitiva utilizando los elementos explicados en la teoría.
\end{enumerate}

\end{document}
